\documentclass[../DoAn.tex]{subfiles}
\begin{document}
\section{Khảo sát hiện trạng}
\label{section:2.1}

Phần này trình bày việc khảo sát về một số hệ thống chatbot tư vấn pháp luật Việt Nam hiện hành trên thị trường, cụ thể bao gồm: AI tra cứu luật\footnote{https://aitracuuluat.vn, truy cập lần cuối ngày 18/05/2026}, AI Luật của LuatVietNam\footnote{https://ailuat.luatvietnam.vn, truy cập lần cuối ngày 18/05/2026}, Chatbot pháp luật của Cổng Pháp luật Quốc gia\footnote{https://ai.phapluat.gov.vn/, truy cập lần cuối ngày 18/05/2026}, Gemini\footnote{https://gemini.google.com, truy cập lần cuối ngày 18/05/2026}.

Hệ thống AITracuuluat là nền tảng trợ lý ảo thông minh cho phép người dùng hỏi đáp pháp luật miễn phí. Bên cạnh việc trả lời câu hỏi pháp luật, công cụ này còn cung cấp các tính năng hỗ trợ chuyên sâu như tra cứu thủ tục hành chính và hệ thống hóa văn bản theo từng lĩnh vực cụ thể. Đây là dự án mang tính xã hội cao, giúp người dân dễ dàng tiếp cận với các quy định pháp luật của Nhà nước mà không cần phải có chuyên môn sâu. Mặc dù vậy, trong quá trình khảo sát câu trả lời của hệ thống với tập các câu hỏi về Luật hôn nhân và gia đình đã phát hiện được một số hạn chế như: Các trích dẫn căn cứ pháp lý đôi khi chưa được cập nhật với các quy định mới nhất và có thể đưa ra các căn cứ pháp lý đã hết hiệu lực. Khi đưa ra các khẳng định trong câu trả lời không trích dẫn được rõ đến từng điều khoản của căn cứ pháp lý mà chỉ tổng hợp lại các căn cứ pháp lý được dùng cho câu trả lời ở cuối mẫu câu trả lời, người dùng phải tự so sánh và đối chiếu. Phần tổng hợp các căn cứ pháp lý ở phía cuối mỗi câu trả lời có thể chứa các căn cứ pháp lý không liên quan đến câu hỏi. Đặc biệt, hệ thống chưa hỗ trợ tính năng cảnh báo hiệu lực văn bản, không có cơ chế đối chiếu để đưa ra các cảnh báo về xung đột pháp lý giữa các quy định khác nhau và cũng chưa có cơ chế để minh bạch hóa luồng xử lý của hệ thống.

AI Luật của LuatVietnam và Chatbot pháp luật của Cổng pháp luật quốc gia là các trợ lý ảo tư vấn pháp luật có khả năng phân tích và đưa ra câu trả lời có độ chính xác cao, có thể trích dẫn chi tiết đến từng điều, khoản của văn bản quy phạm pháp luật đang còn hiệu lực giúp tránh được việc áp dụng các quy định đã cũ hoặc hết hiệu lực. Dẫu vậy, các nền tảng này lại gặp hạn chế lớn về tính đại chúng và khả năng tiếp cận đối với người dân do chỉ miễn phí 5 câu hỏi đầu tiên. Các hệ thống này cũng thiếu tính năng cảnh báo hiệu lực (như thông báo văn bản sắp hết hiệu lực, bị thay thế một phần), chưa có khả năng cảnh báo xung đột pháp lý khi một vấn đề bị chi phối bởi nhiều văn bản luật chồng chéo, và người dùng cũng không thể xem được luồng xử lý của hệ thống.

Gemini là mô hình ngôn ngữ lớn phổ biến toàn cầu, được huấn luyện trên tập dữ liệu phong phú. Tuy nhiên, khi đặt vào ngữ cảnh chuyên biệt của hệ thống pháp luật Việt Nam, mô hình này thường xuyên gặp các vấn đề như: Các câu hỏi phần lớn chưa đưa ra được chính xác và đầy đủ các căn cứ pháp lý; Các căn cứ pháp lý được trích dẫn đôi khi bị hết hiệu lực do dữ liệu huấn luyện hiện nay chỉ đến năm 2025. Vì không có khẳ năng tích hợp cơ sở dữ liệu pháp luật thời gian thực, Gemini hoàn toàn không có khả năng tự động cảnh báo hiệu lực văn bản và xung đột pháp lý cho người dùng.

Tổng hợp lại, bảng \ref{table:fixed_legal_chatbots} dưới đây so sánh tổng quan các chức năng cốt lõi của các chatbot tư vấn pháp luật hiện hành so với định hướng của đồ án.

\begin{table}[H]
\centering
{\small
\renewcommand{\arraystretch}{1.15}
\begin{tabularx}{\textwidth}{|>{\raggedright\arraybackslash}p{2.6cm}|*{4}{>{\centering\arraybackslash}X|}}
\hline
\textbf{Sản phẩm} & \textbf{Tính năng trích dẫn chi tiết căn cứ pháp lý trong câu trả lời} & \textbf{Tính năng cảnh báo hiệu lực} & \textbf{Tính năng cảnh báo xung đột pháp lý} &\textbf{Tính năng hiển thị luồng xử lý của hệ thống}\\ \hline
AITracuuluat & Không & Không & Không & Không \\ \hline
AI Luật & Có & Không & Không & Không \\ \hline
Chatbot Bộ Tư pháp & Có & Không & Không & Không \\ \hline
Gemini & Hạn chế & Không & Không & Có \\ \hline
\end{tabularx}
}
\caption{So sánh các tính năng của các chatbot tư vấn pháp luật hiện hành}
\label{table:fixed_legal_chatbots}
\end{table}

Tóm lại trong các hệ thống chatbot tư vấn pháp luật hiện hành đang tồn tại các vấn đề chính, đó là (i) Đôi khi xảy ra hiện tượng trích dẫn các căn cứ pháp lý thiếu chính xác hoặc không liên quan đến câu hỏi của người dùng. (ii) Chưa có cơ chế cảnh báo tình trạng hiệu lực văn bản (bao gồm văn bản đã hết hiệu lực, sắp hết hiệu lực hoặc sắp có hiệu lực) và cảnh báo về các căn cứ có khả năng mâu thuẫn khi áp dụng trong các tình huống cụ thể. (iii) Chưa có cơ chế hiển thị luồng xử lý, suy luận của chatbot trước khi đưa ra câu trả lời.

Từ những phân tích đã nêu ở trên, ĐATN hướng tới việc xây dựng một hệ thống Hỏi đáp luật hôn nhân gia đình hôn nhân gia đình 2014 có khả năng đưa ra các căn cứ trích dẫn chính xác chi tiết tới từng điều, khoản. Các căn cứ đưa ra được cập nhật mới nhất tới dữ liệu vào tháng 06/2026. Trong các câu trả lời đưa ra có khả năng mô tả ngắn gọn mối quan hệ giữa các căn cứ pháp lý, với mỗi căn cứ pháp lý nêu rõ thời hạn hiệu lực và có thể có các cảnh báo khi tình huống phải áp dụng các văn bản đã hết hiệu lực hoặc sắp bị sửa đổi, thay thế. Đồng thời hệ thống cũng hiển thị chi tiết được luồng xử lý, suy luận trước khi đưa ra câu trả lời. 

\section{Tổng quan chức năng}
\label{section:2.2}
\subsection{Biểu đồ use case tổng quan}
\label{subsection:2.2.1}

Hình \ref{fig:use_case_tong_quan} minh họa biểu đồ use case tổng quan gồm các chức năng mà ĐATN phát triển. Các tác nhân gồm Guest (Khách vãng lai), Registered User (Người dùng đã đăng nhập) và LLM (Mô hình ngôn ngữ lớn). Guest có thể thực hiện: (i) Đăng nhập; (ii) Hỏi đáp luật hôn nhân gia đình; (iii) Xem chi tiết luồng xử lý của các tác nhân; (iv) Xem chi tiết cảnh báo hiệu lực văn bản; (v) Xem chi tiết cảnh báo xung đột pháp lý; (vi) Xem trực quan đồ thị tri thức căn cứ pháp lý. Registered User kế thừa toàn bộ quyền hỏi đáp và kiểm chứng của Guest, đồng thời bổ sung: (i) Đăng xuất; (ii) Quản lý lịch sử hội thoại. LLM đại diện nền tảng trí tuệ nhân tạo bên ngoài, được hệ thống gọi qua API.

\begin{figure}[H]
    \centering
    \includegraphics[width=0.95\textwidth]{Hinhve/Biểu đồ use case tổng quan.png}
    \caption{Biểu đồ use case tổng quan của hệ thống}
    \label{fig:use_case_tong_quan}
\end{figure}
\subsection{Biểu đồ phân rã use case quản lý lịch sử hội thoại}
\label{subsection:2.2.3}

Hình \ref{fig:uc_phan_ra_lich_su} minh họa biểu đồ phân rã use case quản lý lịch sử hội thoại. Nhóm chức năng này dành riêng cho Registered User nhằm quản lý, lưu trữ và tiếp tục các cuộc hội thoại trong quá khứ.
\begin{figure}[H]
    \centering
    \includegraphics[width=0.95\textwidth]{Hinhve/Phân rã use case quản lý lịch sử hội thoại.png}
    \caption{Biểu đồ use case phân rã chức năng Quản lý lịch sử hội thoại}
    \label{fig:uc_phan_ra_lich_su}
\end{figure}

Nhóm chức năng này chỉ dành cho Registered User (người dùng đã đăng nhập vào hệ thống). Cụ thể, Registered User có thể (i) Tạo mới phiên hỏi đáp: khởi tạo và bắt đầu một phiên hỏi đáp mới với chatbot; (ii) Đổi tên phiên hỏi đáp: đổi tên phiên hỏi đáp theo ý muốn (mặc định tên phiên hỏi đáp là câu hỏi đầu tiên của người dùng thuộc phiên này); (iii) Xem lại phiên hỏi đáp cũ: xem lại phiên hỏi đáp cũ của mình, bao gồm cả lịch sử hoạt động của các agent, và tiếp tục hỏi đáp; (iv) Tìm kiếm phiên hỏi đáp: tìm kiếm nhanh phiên hỏi đáp theo tên phiên hỏi đáp; (v) Xóa phiên hỏi đáp: xóa một phiên hỏi đáp bất kỳ khỏi lịch sử hỏi đáp của mình. (vi) Tạo mới project hỏi đáp (vii) Xóa project hỏi đáp (viii) Đổi tên project hỏi đáp (ix) Xem chi tiết project hỏi đáp (x) Di chuyển phiên hỏi đáp vào project hỏi đáp.

\subsection{Quy trình nghiệp vụ Hỏi đáp luật hôn nhân gia đình}
\label{subsection:2.2.4}
\begin{figure}[H]
    \centering
    \includegraphics[width=\textwidth, height=\textheight, keepaspectratio]{Hinhve/Biểu đồ hoạt động quy trình nghiệp vụ Hỏi đáp pháp luật.png}
    \caption{Biểu đồ hoạt động quy trình nghiệp vụ Hỏi đáp luật hôn nhân gia đình}
    \label{fig:activity_business_process}
\end{figure}

Hình \ref{fig:activity_business_process} mô tả chi tiết quy trình nghiệp vụ use case Hỏi đáp luật hôn nhân gia đình như sau: (i) Người dùng truy cập vào ứng dụng, đăng nhập bằng Google OAuth, hệ thống xác thực tài khoản, tải lại lịch sử hội thoại từ PostgreSQL và gắn phiên hỏi đáp với tài khoản; (ii) Người dùng nhập câu hỏi liên quan đến luật hôn nhân gia đình bằng ngôn ngữ tự nhiên. (iii) Router Agent phân tích câu hỏi và đề xuất danh sách retriever cần gọi. Danh sách này không được thực thi ngay mà phải đi qua bước kiểm tra tính hợp lệ, tính trùng lặp và retriever policy để loại bỏ retriever không tồn tại, retriever trùng lặp (bị gọi nhiều lần) hoặc retriever không phù hợp (iv) Các retriever hợp lệ lựa chọn template Cypher query, trích xuất tham số có cấu trúc và truy vấn đồ thị tri thức để lấy context pháp lý. (v) Context được lọc theo hiệu lực, loại bỏ trùng lặp, phân loại căn cứ chính, căn cứ bổ trợ, căn cứ hướng dẫn, căn cứ thay thế và gán các cờ cảnh báo về hiệu lực hoặc mâu thuẫn. (vi) Trước khi gọi API LLM để sinh câu trả lời, hệ thống kiểm tra context đã chuẩn hóa có căn cứ pháp lý để sinh câu trả lời hay không. Nếu context không có căn cứ pháp lý, hệ thống cần thông báo không trả lời được, yêu cầu người dùng đặt câu hỏi khác hoặc bổ sung thông tin. (vii) Khi có căn cứ trong context, hệ thống kết hợp câu hỏi, lịch sử hội thoại và context đã chuẩn hóa vào prompt để gọi API LLM sinh câu trả lời. (viii) Hệ thống hiển thị câu trả lời, Agent Trace, đường dẫn trực quan hóa phần đồ thị tri thức đã truy xuất, các cảnh báo hiệu lực và (hoặc) mâu thuẫn để người dùng hoặc Provider có thể đối chiếu.

\section{Đặc tả chức năng}
\subsection{Đặc tả use case Hỏi đáp luật hôn nhân gia đình}
Bảng \ref{table:uc_specs_qa} là bảng đặc tả use case Hỏi đáp luật hôn nhân gia đình.
\begin{table}[H]
\centering
\begin{tabular}{|p{3.5cm}|p{11.5cm}|}
\hline
\textbf{Mã usecase:} UC01 & \textbf{Tên usecase:} Hỏi đáp luật hôn nhân gia đình\\ \hline
\textbf{Tác nhân} & Guest (Khách vãng lai) hoặc Registered User (Người dùng hệ thống), LLM (Mô hình ngôn ngữ lớn được gọi thông qua API) \\ \hline
\textbf{Mô tả vắn tắt} & Tiếp nhận câu hỏi của người dùng và sinh câu trả lời \\ \hline
\textbf{Tiền điều kiện} & Người dùng đã truy cập vào giao diện hội thoại chính của ứng dụng Chatbot. \\ \hline
\textbf{Luồng chính} & 
1. Người dùng nhập nội dung câu hỏi pháp lý vào ô nhập liệu và nhấn nút gửi. \newline
2. Hệ thống tiếp nhận câu hỏi, hiển thị tin nhắn lên màn hình trò chuyện và khởi tạo giao diện hiển thị luồng phân tích của tác nhân. \newline
3. Hệ thống chuyển giao câu hỏi và lịch sử hội thoại sang cho Router Agent. \newline
4. Hệ thống thực hiện phân tích ý định câu hỏi thông qua Tác nhân định tuyến nhằm xác định các Retrieval Agents cần gọi. \newline
5. Hệ thống nhận về context pháp lý truy xuất từ đồ thị tri thức \newline
6. Hệ thống thực hiện chuẩn hóa ngữ cảnh thu được, gán các cờ cảnh báo về hiệu lực, sửa đổi bổ sung, mâu thuẫn. \newline
7. Hệ thống kết hợp ngữ cảnh đã được chuẩn hóa vào system prompt rồi gửi yêu cầu đến API LLM. \newline
8. LLM xử lý dữ liệu và sinh câu trả lời lập luận tư vấn pháp lý hoàn chỉnh theo đúng template quy định. \newline
9. Hệ thống đóng giao diện tiến trình trung gian và hiển thị câu trả lời hoàn chỉnh của chatbot lên màn hình người dùng. \\ \hline
\textbf{Luồng thay thế 1} & 
4.1. Hệ thống nhận được danh sách các retriever agent cần gọi rỗng do Router Agent phân tích ý định nhưng không lựa chọn được retriever chuyên biệt phù hợp hoặc lựa chọn retriever không tồn tại. \newline
4.2. Hệ thống tự động kích hoạt retriever dự phòng mặc định \newline
4.3. Quay lại bước 5 của luồng chính\\ \hline
\textbf{Luồng thay thế 2} & 
2.1. Người dùng nhấn nút dừng tiến trình ("Stop") trên giao diện ô chat khi chatbot đang trong quá trình suy luận hoặc sinh text. \newline
2.2. Hệ thống ngay lập tức ngắt kết nối API và dừng việc hiển thị dữ liệu. use case kết thúc tại đây. \\ \hline
\textbf{Hậu điều kiện} & Hệ thống hiển thị thành công câu trả lời tư vấn pháp lý kèm các căn cứ trích dẫn, khối Agent Trace ở trạng thái thu gọn, các nút báo hiệu lực và xung đột pháp lý và đường dẫn liên kết truy cập đồ thị tri thức căn cứ pháp lý phía dưới mỗi câu trả lời.\\ \hline
\end{tabular}
\caption{Đặc tả chi tiết use case Hỏi đáp luật hôn nhân gia đình}
\label{table:uc_specs_qa}
\end{table}
\subsection{Đặc tả use case Xem chi tiết luồng xử lý của các tác nhân}
Bảng \ref{table:uc_specs_trace} đặc tả chi tiết cho use case Xem chi tiết luồng xử lý của các tác nhân (Agent Trace). 
\begin{table}[H]
\centering
\begin{tabular}{|p{3.5cm}|p{11.5cm}|}
\hline
\textbf{Mã usecase:} UC02 & \textbf{Tên usecase:} Xem chi tiết luồng xử lý của các tác nhân \\ \hline
\textbf{Tác nhân} & Guest (Khách vãng lai) hoặc Registered User (Người dùng hệ thống) \\ \hline
\textbf{Mô tả vắn tắt} & Cho phép người dùng chủ động mở rộng giao diện tiến trình trung gian để kiểm tra chi tiết các agent được kích hoạt và dữ liệu ngữ cảnh thô trích xuất từ đồ thị tri thức. \\ \hline
\textbf{Tiền điều kiện} & Người dùng đã thực hiện gửi câu hỏi và chatbot đã hoàn tất việc trả lời câu hỏi. \\ \hline
\textbf{Luồng chính} & 
1. Tại giao diện tin nhắn của phiên hội thoại hiện tại, người dùng di chuyển đến vùng hiển thị tiến trình của các tác nhân (Agent Trace). \newline
2. Người dùng nhấp chuột vào thanh tiêu đề của khối "Luồng truy xuất ngữ cảnh" hoặc "Tổng hợp đáp án". Giao diện mở khối nội dung theo cấu trúc phân cấp.\newline
3. Hệ thống hiển thị trạng thái đang gọi tác nhân Router Agent, sau đó lần lượt là các retrievers được gọi. \newline
4. Người dùng cuộn chuột xuống khối "Tổng hợp đáp án", nhấp chuột vào đó để kiểm tra hộp mã nguồn định dạng JSON, chứa toàn bộ các thuộc tính và nội dung của các nút văn bản pháp luật trích xuất từ Neo4j và đã được chuẩn hóa thành context cho LLM. \newline
5. Người dùng nhấp lại vào thanh tiêu đề một lần nữa để thu gọn khối hiển thị sau khi hoàn tất việc kiểm chứng. \\ \hline
\textbf{Luồng thay thế 1} & 
3.1. Người dùng yêu cầu mở rộng tiến trình nhưng có lỗi từ hệ thống. \newline
3.2. Hệ thống hiển thị thông báo lỗi: "Không tìm thấy dữ liệu luồng xử lý của phiên hội thoại này". \newline
3.3. Người dùng nhấp đóng thông báo lỗi. use case kết thúc tại đây. \\ \hline
\textbf{Hậu điều kiện} & Hệ thống hiển thị chính xác, đầy đủ và trực quan toàn bộ các tham số đầu vào/đầu ra của các agent.  \\ \hline
\end{tabular}
\caption{Đặc tả chi tiết use case Xem chi tiết luồng xử lý của các tác nhân}
\label{table:uc_specs_trace}
\end{table}
\subsection{Đặc tả use case Xem chi tiết cảnh báo hiệu lực văn bản}
Bảng \ref{table:uc_specs_validity} đặc tả chi tiết use case Xem chi tiết cảnh báo hiệu lực văn bản.
\begin{table}[H]
\centering
\begin{tabular}{|p{3.5cm}|p{11.5cm}|}
\hline
\textbf{Mã usecase:} UC03 & \textbf{Tên usecase:} Xem chi tiết cảnh báo hiệu lực văn bản \\ \hline
\textbf{Tác nhân} & Guest (Khách vãng lai) hoặc Registered User (Người dùng hệ thống) \\ \hline
\textbf{Mô tả vắn tắt} & Cho phép người dùng xem chi tiết về tình trạng hiệu lực của văn bản pháp luật được trích dẫn trong câu trả lời. \\ \hline
\textbf{Tiền điều kiện} & Trong phiên hội thoại hiện tại, chatbot đã hoàn tất sinh ít nhất một câu trả lời có chứa ít nhất một cảnh báo hiệu lực văn bản. \\ \hline
\textbf{Luồng chính} & 
1. Người dùng nhấp chuột vào biểu tượng cảnh báo hiệu lực và mâu thuẫn ở phía cuối câu trả lời. \newline
2. Hệ thống hiển thị Panel chi tiết bao gồm tên các căn cứ pháp lý, trạng thái hiệu lực (còn hiệu lực, sắp hết hiệu lực, đã hết hiệu lực hoặc bị thay thế) kèm thời gian hiệu lực, hết hiệu lực, và căn cứ thay thế (nếu có). \newline
3. Người dùng đọc thông tin, có thể nhấp liên kết tới điều khoản căn cứ để xem nộ dung chi tiết căn cứ ở trang thuvienphapluat. \newline
4. Người dùng đóng panel sau khi hoàn tất việc đối chiếu. \\ \hline
\textbf{Hậu điều kiện} & Người dùng nắm được chi tiết tình trạng hiệu lực của các căn cứ pháp lý trong câu trả lời. \\ \hline
\end{tabular}
\caption{Đặc tả chi tiết use case Xem chi tiết cảnh báo hiệu lực văn bản}
\label{table:uc_specs_validity}
\end{table}
\subsection{Đặc tả use case Xem chi tiết cảnh báo xung đột pháp lý}
Bảng \ref{table:uc_specs_conflict} đặc tả chi tiết use case Xem chi tiết cảnh báo xung đột pháp lý.
\begin{table}[H]
\centering
\begin{tabular}{|p{3.5cm}|p{11.5cm}|}
\hline
\textbf{Mã usecase:} UC04 & \textbf{Tên usecase:} Xem chi tiết cảnh báo xung đột pháp lý \\ \hline
\textbf{Tác nhân} & Guest (Khách vãng lai) hoặc Registered User (Người dùng hệ thống) \\ \hline
\textbf{Mô tả vắn tắt} & Cho phép người dùng xem chi tiết về mâu thuẫn hoặc chồng chéo giữa các căn cứ pháp lý trong câu trả lời. \\ \hline
\textbf{Tiền điều kiện} & Trong phiên hội thoại hiện tại, chatbot đã hoàn tất sinh ít nhất một câu trả lời có chứa ít nhất một cảnh báo xung đột pháp lý. \\ \hline
\textbf{Luồng chính} & 
1. Người dùng nhấp chuột vào biểu tượng cảnh báo hiệu lực và mâu thuẫn ở phía cuối câu trả lời. \newline
2. Hệ thống hiển thị Panel chi tiết bao gồm tên các căn cứ pháp lý có mâu thuẫn, mô tả ngắn về nội dung mâu thuẫn. \newline
3. Người dùng đọc thông tin, có thể nhấp liên kết tới từng điều khoản liên quan để xem chi tiết tại trang thuvienphapluat \newline
4. Người dùng đóng panel sau khi hoàn tất việc đối chiếu. \\ \hline
\textbf{Hậu điều kiện} & Người dùng nắm được chi tiết các căn cứ pháp lý có thể gây ra xung đột. \\ \hline
\end{tabular}
\caption{Đặc tả chi tiết use case Xem chi tiết cảnh báo xung đột pháp lý}
\label{table:uc_specs_conflict}
\end{table}
\subsection{Đặc tả use case Xem trực quan đồ thị căn cứ pháp lý của câu trả lời}
Bảng \ref{table:uc_specs_kg} đặc tả chi tiết use case Xem trực quan đồ thị căn cứ pháp lý của câu trả lời.
\begin{table}[H]
\centering
\begin{tabular}{|p{3.5cm}|p{11.5cm}|}
\hline
\textbf{Mã usecase:} UC05 & \textbf{Tên usecase:} Xem trực quan đồ thị căn cứ pháp lý của câu trả lời \\ \hline
\textbf{Tác nhân} & Guest (Khách vãng lai) hoặc Registered User (Người dùng hệ thống) \\ \hline
\textbf{Mô tả vắn tắt} & Cho phép người dùng mở giao diện trực quan hoá phần đồ thị tri thức đã được truy xuất cho câu trả lời. \\ \hline
\textbf{Tiền điều kiện} & Trong phiên hội thoại hiện tại, chatbot đã hoàn tất sinh ít nhất một câu trả lời. \\ \hline
\textbf{Luồng chính} & 
1. Tại câu trả lời của chatbot, người dùng nhấp chuột vào liên kết ``Xem đồ thị căn cứ pháp lý'' ở cuối câu trả lời \newline
2. Hệ thống tải dữ liệu đồ thị tri thức tương ứng với lượt hỏi đáp và mở giao diện trực quan hoá trong 1 tab mới. \newline
3. Giao diện hiển thị các nút văn bản, điều khoản và quan hệ giữa chúng đã được truy xuất. \newline
4. Người dùng tương tác với đồ thị. \newline
5. Người dùng đóng tab sau khi hoàn tất việc kiểm chứng. \\ \hline
\textbf{Luồng thay thế 1} & 
2.1. Hệ thống không tải được dữ liệu đồ thị tri thức. \newline
2.2. Giao diện hiển thị thông báo: ``Không tìm thấy dữ liệu đồ thị cho câu trả lời này''. \newline
2.3. Người dùng đóng thông báo. use case kết thúc tại đây. \\ \hline
\textbf{Hậu điều kiện} & Người dùng quan sát được phần đồ thị tri thức được truy xuất để trả lời câu hỏi. \\ \hline
\end{tabular}
\caption{Đặc tả chi tiết use case Xem trực quan đồ thị tri thức căn cứ pháp lý}
\label{table:uc_specs_kg}
\end{table}
\subsection{Đặc tả use case Xem lại phiên hỏi đáp cũ}
Bảng \ref{table:uc_specs_resume} đặc tả use case Xem lại phiên hỏi đáp cũ thuộc use case Quản lý lịch sử hội thoại.
\begin{table}[H]
\centering
\begin{tabular}{|p{3.5cm}|p{11.5cm}|}
\hline
\textbf{Mã usecase:} UC06 & \textbf{Tên usecase:} Xem lại phiên hỏi đáp cũ \\ \hline
\textbf{Tác nhân} & Registered User (Người dùng hệ thống) \\ \hline
\textbf{Mô tả vắn tắt} & Cho phép người dùng đã đăng nhập truy cập vào danh mục lịch sử trò chuyện ở thanh điều hướng để khôi phục toàn bộ nội dung tin nhắn, agent trace, đường dẫn đồ thị tri thức căn cứ pháp lý và cảnh báo hiệu lực, mâu thuẫn. \\ \hline
\textbf{Tiền điều kiện} & Người dùng đã đăng nhập thành công vào hệ thống và có ít nhất một phiên hội thoại cũ. \\ \hline
\textbf{Luồng chính} & 
1. Người dùng di chuyển chuột đến thanh điều hướng của giao diện Chainlit để xem danh sách các phiên hội thoại cũ. \newline
2. Giao diện hiển thị danh sách tiêu đề của các phiên hội thoại quá khứ theo thứ tự thời gian tương tác gần nhất hoặc theo các project hỏi đáp \newline
3. Người dùng nhấp chuột lựa chọn vào một phiên hội thoại cụ thể muốn xem lại. \newline
4. Hệ thống tiếp nhận sự kiện, tải toàn bộ danh sách tin nhắn, agent trace, đường dẫn đồ thị tri thức căn cứ pháp lý và cảnh báo hiệu lực, mâu thuẫn \newline
5. Hệ thống khôi phục và hiển thị chính xác toàn bộ lịch sử trò chuyện. \\ \hline
\textbf{Hậu điều kiện} & Hệ thống hiển thị thành công phiên hội thoại cũ cùng các agent trace, đường dẫn đồ thị tri thức và cảnh báo hiệu lực, mâu thuẫn. Cửa sổ chat sẵn sàng tiếp nhận các câu hỏi tiếp theo của người dùng cho phiên hội thoại này. \\ \hline
\end{tabular}
\caption{Đặc tả chi tiết use case Xem lại phiên hỏi đáp cũ}
\label{table:uc_specs_resume}
\end{table}
\section{Yêu cầu phi chức năng}
\label{section:2.4}
\subsection{Yêu cầu chung về hệ thống}
Thời gian phản hồi trung bình của hệ thống đối với mỗi truy vấn của người dùng phải đạt mức dưới 30 giây.
\subsection{Yêu cầu về giao diện người dùng (UI/UX)}
\label{subsection:2.4.2}
Giao diện xây dựng trên nền tảng Chainlit phải có bố cục phân tách rõ ràng giữa khu vực Sidebar lưu lịch sử và khung trò chuyện chính đối với Registered User; với Guest, thanh điều hướng lịch sử hội thoại không được hiển thị. Khối hiển thị luồng hoạt động trung gian của các tác nhân phải được thiết kế dưới dạng cấu trúc cây phân cấp có thể đóng/mở linh hoạt. Dữ liệu trích xuất từ các retrievers phải có định dạng rõ ràng, tường minh để người dùng dễ dàng đối chiếu, kiểm chứng căn cứ pháp lý. Các thông báo từ hệ thống nằm ở góc trên bên phải. Giao diện mặc định có 2 theme Dark và Light. Giao diện tab trực quan hóa đồ thị tri thức phải có khả năng phóng to, thu nhỏ, di chuyển các nút và hiển thị id của từng nút khi người dùng di chuột vào; đồng thời có màu sắc khác nhau cho các loại node khác nhau (căn cứ chính, căn cứ hướng dẫn, căn cứ bổ trợ, căn cứ thay thế, căn cứ mâu thuẫn).
\subsection{Yêu cầu về giải pháp kỹ thuật đề xuất (Agentic GraphRAG)}
\label{subsection:2.4.3}
Để đánh giá hiệu năng và chất lượng đầu ra của hệ thống, các chỉ số mục tiêu được đặt ra như sau: (i) \textbf{Nhóm Metric RAGAS} bao gồm Faithfulness, Context Recall đạt tối thiểu 85\%; (ii) \textbf{Nhóm căn cứ pháp lý} bao gồm Legal Validity Accuracy đạt tối thiểu 90\% và Legal Citation Accuracy tối thiểu 85\%; (iii) \textbf{Nhóm xử lý xung đột} bao gồm Legal Conflict Handling Accuracy tối thiểu 80\% và Context Conflict Detection Rate đạt tối thiểu 80\%; (iv) \textbf{Nhóm phân loại viện dẫn} bao gồm Citation Prioritization Accuracy đạt tối thiểu 85\% và Invalid Citation Rate dưới 10\%. Chi tiét các metric được trình bày trong chương \ref{chapter:Experiment}.
%%%%%%%%%%%%%%%%%%%%%%%%%%%%%%%%%%%
\end{document}
